\documentclass[12pt,a4paper]{article}

\usepackage{fontspec}
\usepackage{polyglossia}
\usepackage{geometry}
\usepackage{array}
\usepackage{longtable}
\usepackage{booktabs}
\usepackage{enumitem}
\usepackage{setspace}
\usepackage{hyperref}

\geometry{margin=2.2cm}
\setstretch{1.2}
\setmainlanguage{russian}
\setotherlanguage{english}
\setmainfont{Times New Roman}
\setsansfont{Arial}
\setmonofont{Courier New}
\newfontfamily\cyrillicfont{Times New Roman}
\newfontfamily\cyrillicfontsf{Arial}
\newfontfamily\cyrillicfonttt{Times New Roman}
\hypersetup{
    colorlinks=true,
    linkcolor=black,
    urlcolor=blue,
    citecolor=black
}
\urlstyle{same}

\begin{document}

\begin{titlepage}
    \begin{center}
        \textbf{Университет ИТМО}\\
        Факультет программной инженерии и компьютерной техники\\[2cm]

        \textbf{\Large Отчёт по групповому проекту}\\[0.3cm]
        \textbf{\large по курсу <<Разработка мобильных приложений>>}\\[1.4cm]

        \textbf{\Large Система отслеживания акций и портфеля}\\[0.4cm]
        Android-клиент, backend, сервис котировок и observability-контур\\[2cm]
    \end{center}

    \begin{flushright}
        \begin{tabular}{rl}
            Студент: & \underline{\hspace{6.5cm}} \\
            Группа: & \underline{\hspace{6.5cm}} \\
            Команда проекта: & \underline{\hspace{6.5cm}} \\
            Преподаватель: & Ключев А.\,О. \\
        \end{tabular}
    \end{flushright}

    \vfill

    \begin{center}
        Санкт-Петербург\\
        2026
    \end{center}
\end{titlepage}

\tableofcontents
\newpage

\section*{Введение}
\addcontentsline{toc}{section}{Введение}

Групповой проект по курсу <<Разработка мобильных приложений>> направлен на создание многокомпонентной системы для просмотра рыночных котировок, ведения пользовательского портфеля и выполнения операций покупки и продажи акций. На момент подготовки этого отчёта проект уже вышел за пределы монолитного учебного backend и включает Android-клиент, мобильный API-шлюз, транзакционный backend, отдельный сервис котировок на Go, а также контейнеризированный стек observability.

Отчёт подготовлен по актуальному состоянию репозитория после обновления ветки \texttt{main} на 1 апреля 2026 года и повторного импорта Android-клиента из внешнего репозитория команды. Старые feature-ветки общего репозитория рассмотрены как этапы развития backend-части, а архитектурный baseline определяется именно актуальным \texttt{main}. Структура отчёта соответствует требованиям лекционных материалов курса: техническое задание, описание архитектуры, описание реализации и описание системы тестирования~[1].

\section{Техническое задание}

\subsection{Введение}

Разрабатываемая система предназначена для демонстрации архитектурных и прикладных аспектов мобильной разработки: пользовательский интерфейс Android-приложения, взаимодействие с backend API, работа с котировками, хранение транзакционных и аналитических данных, а также сбор telemetry-данных для диагностики распределённой системы.

Проект рассматривается как учебная торгово-информационная платформа. Пользователь системы должен иметь возможность зарегистрироваться, войти в приложение, получить актуальную котировку по тикеру, выполнить покупку или продажу актива и увидеть агрегированную статистику своего портфеля.

\subsection{Назначение разработки}

Назначение разработки заключается в создании многомодульного мобильного программного комплекса, который:

\begin{itemize}[leftmargin=1.2cm]
    \item предоставляет Android-клиент для взаимодействия с пользователем;
    \item предоставляет backend API для аутентификации и операций с портфелем;
    \item получает и публикует рыночные котировки;
    \item демонстрирует работу распределённой архитектуры с observability-подсистемой;
    \item обеспечивает базу для дальнейшего расширения системы тестирования и документации.
\end{itemize}

\subsection{Требования к программе или программному изделию}

\subsubsection{Функции системы}

Система должна обеспечивать следующие функции:

\begin{enumerate}[leftmargin=1.2cm]
    \item Регистрация пользователя по логину и паролю.
    \item Аутентификация пользователя и выдача access token.
    \item Просмотр актуальной котировки по тикеру.
    \item Просмотр информации о владении выбранной акцией в портфеле.
    \item Выполнение операций покупки и продажи акций.
    \item Получение агрегированной статистики портфеля.
    \item Ведение истории котировок и публикация событий в инфраструктурные каналы.
    \item Сбор трассировок и технических данных для анализа работы сервисов.
\end{enumerate}

\subsubsection{Состав системы}

В состав системы входят:

\begin{itemize}[leftmargin=1.2cm]
    \item Android-клиент \texttt{android-app/};
    \item публичный API gateway \texttt{mobile-api/};
    \item транзакционный сервис портфеля \texttt{mobile-backend/};
    \item сервис котировок \texttt{quotes-service/};
    \item Linux C-драйвер \texttt{driver/} как низкоуровневый источник котировок;
    \item PostgreSQL для транзакционных данных;
    \item Redis для событий и latest-cache;
    \item ClickHouse для истории котировок;
    \item OpenTelemetry Collector, Tempo, Prometheus, Loki и Grafana для наблюдаемости;
    \item документация проекта и база знаний в Obsidian.
\end{itemize}

\subsubsection{Характеристики системы}

К системе предъявляются следующие характеристики:

\begin{itemize}[leftmargin=1.2cm]
    \item модульность и разделение ответственности между клиентом, gateway, backend и сервисом котировок;
    \item возможность локального запуска в Docker Compose;
    \item расширяемость по API, хранилищам и observability;
    \item поддержка трассировки запросов через OpenTelemetry;
    \item наличие автоматизированных тестов хотя бы для части backend-компонентов;
    \item документированность архитектуры, реализации и тестирования;
    \item возможность дальнейшей интеграции Android-клиента с реальным backend API.
\end{itemize}

\section{Описание архитектуры}

\subsection{Архитектурная идея}

Ключевая идея проекта состоит в разделении системы на несколько относительно независимых подсистем. Такое разделение позволяет не смешивать клиентский UI, бизнес-логику портфеля, рыночные данные, техническую инфраструктуру и observability. Для учебного проекта это особенно важно, так как даёт возможность параллельной работы нескольких участников команды и упрощает понимание внутренних границ системы.

В актуальном состоянии архитектура имеет сервисный характер:

\begin{itemize}[leftmargin=1.2cm]
    \item Android-клиент отвечает за пользовательский интерфейс;
    \item \texttt{mobile-api} является внешней точкой входа для мобильных приложений;
    \item \texttt{mobile-backend} реализует бизнес-логику аутентификации, портфеля и торговых операций;
    \item \texttt{quotes-service} поставляет актуальные котировки и историю;
    \item инфраструктурный слой хранит данные и собирает телеметрию.
\end{itemize}

\subsection{Декомпозиция на компоненты}

\begin{longtable}{p{0.23\textwidth} p{0.7\textwidth}}
\toprule
\textbf{Компонент} & \textbf{Архитектурная роль} \\
\midrule
\endhead
\texttt{android-app} & Android-прототип с экранами логина, профиля, статистики, списка акций и карточки акции. Пока интегрирован с backend лишь концептуально. \\
\texttt{mobile-api} & Ktor gateway, который принимает клиентские HTTP-запросы, выполняет общие middleware-функции и проксирует запросы во внутренние сервисы. \\
\texttt{mobile-backend} & Транзакционный сервис портфеля. Управляет пользователями, JWT, операциями покупки и продажи, статистикой и взаимодействием с PostgreSQL. \\
\texttt{quotes-service} & Go-сервис чтения котировок, поддерживающий mock-источник и Linux-драйвер, а также публикацию данных в Redis и ClickHouse. \\
\texttt{driver} & Симуляция устройства \texttt{/dev/quotes} на стороне Linux-хоста. Нужен как низкоуровневая часть учебной задачи по источнику рыночных данных. \\
\texttt{postgres} & Основное хранилище пользователей, портфелей, lot-структур и истории сделок. \\
\texttt{redis} & Канал Streams для событий и быстрый storage latest-cache по котировкам. \\
\texttt{clickhouse} & Хранилище истории котировок, удобное для аналитических и временных запросов. \\
\texttt{otel-collector} и related stack & Приём и маршрутизация telemetry-сигналов от сервисов, интеграция с Tempo, Prometheus, Loki и Grafana. \\
\bottomrule
\end{longtable}

\subsection{Внутренняя архитектура Kotlin backend}

Сервис \texttt{mobile-backend} построен по модели слоистой архитектуры, близкой к Clean Architecture~[2]. Внутренние пакеты разделены на:

\begin{itemize}[leftmargin=1.2cm]
    \item \texttt{presentation} --- HTTP-маршруты, DTO, плагины Ktor, обработка ошибок;
    \item \texttt{application} --- use case и порты;
    \item \texttt{domain} --- предметные сущности и value objects;
    \item \texttt{infrastructure} --- БД, безопасность, Redis, OpenTelemetry, HTTP-клиенты и конфигурация.
\end{itemize}

Такое разделение делает use case слабее связанными с Ktor, JDBC и конкретными внешними библиотеками. В результате тестирование логики регистрации, торговли и статистики возможно без полного запуска серверного окружения.

\subsection{Основные потоки взаимодействия}

\subsubsection{Сценарий аутентификации и торговых операций}

\begin{enumerate}[leftmargin=1.2cm]
    \item Клиент отправляет запрос в \texttt{mobile-api}.
    \item \texttt{mobile-api} создаёт trace/span, фиксирует request context и пересылает запрос в \texttt{mobile-backend}.
    \item \texttt{mobile-backend} выполняет use case, при необходимости обращается к PostgreSQL и публикует событие в Redis Streams.
    \item Ответ возвращается назад через \texttt{mobile-api}.
\end{enumerate}

\subsubsection{Сценарий получения котировок}

\begin{enumerate}[leftmargin=1.2cm]
    \item \texttt{quotes-service} регулярно считывает снапшоты котировок из mock-файла или из \texttt{/dev/quotes}.
    \item Сервис обновляет in-memory state, latest-cache в Redis и историю в ClickHouse.
    \item \texttt{mobile-backend} обращается к \texttt{quotes-service} за актуальной ценой.
    \item \texttt{mobile-api} возвращает котировку мобильному клиенту.
\end{enumerate}

\subsection{Хранилища и интеграции}

Хранилища выбраны не случайно, а по ролям:

\begin{itemize}[leftmargin=1.2cm]
    \item PostgreSQL хранит транзакционные данные, где важны консистентность и работа с отношениями;
    \item Redis используется для Streams и быстрого latest-cache;
    \item ClickHouse хранит историю котировок как аналитический временной ряд.
\end{itemize}

Такое разделение упрощает дальнейший рост проекта: транзакционные и аналитические нагрузки не смешиваются, а события можно использовать для новых фоновых сервисов.

\subsection{Наблюдаемость и OpenTelemetry}

В проекте уже реализован базовый контур наблюдаемости:

\begin{itemize}[leftmargin=1.2cm]
    \item \texttt{mobile-api} создаёт серверные span на HTTP-запросы;
    \item \texttt{mobile-backend} использует \texttt{OpenTelemetryTelemetryRecorder} для внутренних span;
    \item \texttt{quotes-service} подключает OTLP exporter для Go;
    \item \texttt{otel-collector} принимает traces, metrics и logs и передаёт traces в Tempo~[6].
\end{itemize}

Архитектурно это важно по двум причинам. Во-первых, observability встроена в систему не как внешняя надстройка, а как часть её runtime-модели. Во-вторых, наличие gateway и нескольких backend-компонентов делает распределённую трассировку действительно полезной.

\subsection{Анализ веток и источников проекта}

При подготовке архитектурного описания важно учитывать историю репозитория. В проекте присутствуют ветки \texttt{origin/feature/market-data} и \texttt{origin/feature/mobile-backend-db}, однако обе они уже не содержат уникальных коммитов относительно актуального \texttt{origin/main}. Следовательно, они представляют ранние этапы развития backend, но не определяют текущее состояние системы.

Android-клиент, наоборот, не находился в общем Git-источнике проекта и был импортирован отдельно из внешнего репозитория команды. Поэтому архитектурно проект сейчас состоит из двух линий происхождения: основной backend-контур из общего репозитория и отдельно созданный клиентский Android-прототип. Это необходимо честно фиксировать в отчёте, чтобы не возникало ложного впечатления, будто весь проект развивался в одном и том же дереве исходников с самого начала.

\subsection{Архитектурные ограничения}

Несмотря на сильную серверную часть, есть и важные ограничения:

\begin{itemize}[leftmargin=1.2cm]
    \item Android-клиент пока не работает с реальным API проекта;
    \item в карточке акции покупка и продажа ещё не привязаны к backend;
    \item часть UI-статистики и графиков работает на моках;
    \item у \texttt{mobile-api} пока нет собственного тестового набора.
\end{itemize}

Эти ограничения не отменяют архитектуру, но показывают, где проходит граница между уже реализованным контуром и следующим этапом интеграции.

\section{Описание реализации}

\subsection{Android-клиент}

Импортированный Android-клиент реализован на Kotlin с использованием классического стека Android Views, XML layouts, \texttt{Activity}, \texttt{Fragment}, \texttt{RecyclerView} и \texttt{BottomNavigationView}~[4]. Основные точки входа:

\begin{itemize}[leftmargin=1.2cm]
    \item \texttt{MainActivity} --- экран входа;
    \item \texttt{SecondActivity} --- контейнер нижней навигации;
    \item \texttt{ProfileFragment}, \texttt{StatisticsFragment}, \texttt{StockListFragment}, \texttt{StockChartFragment} --- основные пользовательские экраны.
\end{itemize}

Текущая реализация клиента носит прототипный характер. Вход реализован локальной проверкой непустых полей, котировки берутся напрямую из MOEX ISS через \texttt{HttpURLConnection}, а в части экранов используются локально сгенерированные значения. Тем не менее экранная структура уже хорошо подготовлена для подключения реального backend API.

\subsection{API gateway}

Модуль \texttt{mobile-api} реализован на Ktor~[3]. Его \texttt{ApplicationModule} выполняет:

\begin{itemize}[leftmargin=1.2cm]
    \item загрузку конфигурации;
    \item создание OpenTelemetry handle;
    \item создание HTTP-клиента для upstream;
    \item подключение Ktor plugins: call id, logging, content negotiation, tracing, status pages, routing.
\end{itemize}

Маршруты gateway повторяют клиентский API-контракт и проксируют вызовы к \texttt{mobile-backend}. Такой слой полезен не только как reverse proxy, но и как архитектурный буфер между мобильным приложением и внутренними сервисами.

\subsection{Транзакционный backend}

Модуль \texttt{mobile-backend} является самым зрелым функциональным слоем проекта. В bootstrap создаются:

\begin{itemize}[leftmargin=1.2cm]
    \item use case аутентификации;
    \item use case торговли;
    \item use case чтения деталей владения;
    \item use case получения статистики;
    \item репозитории Exposed;
    \item publisher событий в Redis Streams;
    \item компонент для записи telemetry-событий;
    \item JWT и password hashing инфраструктура.
\end{itemize}

Из реализованных HTTP-операций в backend доступны регистрация, вход, получение котировки, получение информации по акции, покупка, продажа и статистика портфеля. На текущем этапе это ядро серверной части всей системы.

\subsection{Сервис котировок и инфраструктура}

\texttt{quotes-service} реализован на Go и запускает HTTP-сервер с маршрутом \texttt{/health}, \texttt{/metrics}, \texttt{/quotes}, \texttt{/quotes/\{ticker\}} и \texttt{/quotes/\{ticker\}/history}. Сервис умеет:

\begin{itemize}[leftmargin=1.2cm]
    \item читать котировки из mock-файла или из Linux-источника;
    \item держать актуальное состояние в памяти;
    \item публиковать latest-state и события в Redis;
    \item сохранять историю в ClickHouse;
    \item оборачивать HTTP-обработчики через OpenTelemetry middleware.
\end{itemize}

Верхний уровень инфраструктуры описан в \texttt{docker-compose.yml}. Он поднимает \texttt{caddy}, \texttt{postgres}, \texttt{redis}, \texttt{clickhouse}, \texttt{quotes-service}, \texttt{portfolio-service}, \texttt{mobile-api}, \texttt{otel-collector}, \texttt{loki}, \texttt{tempo}, \texttt{prometheus} и \texttt{grafana}. Таким образом, проект уже имеет полноценный локальный стенд для разработки и демонстрации.

\subsection{Документация и база знаний}

В рамках подготовки сдачи были заново сформированы:

\begin{itemize}[leftmargin=1.2cm]
    \item проектная документация в каталоге \texttt{docs/};
    \item база знаний Obsidian в каталоге \texttt{knowledge-base/obsidian/};
    \item исходный LaTeX-текст отчёта и PDF-сборка.
\end{itemize}

Такой набор материалов соответствует требованиям курса: знания фиксируются не только в коде, но и в поясняющих артефактах, которые можно использовать при защите и передаче проекта внутри команды.

\section{Описание системы тестирования}

\subsection{Требуемые уровни тестирования}

Согласно лекционным материалам, в отчёте должны быть отражены три уровня тестирования: модульное, интеграционное и системное~[1]. В текущем проекте это особенно важно, так как система распределена по нескольким технологиям и языкам.

\subsection{Текущее покрытие тестами}

\subsubsection{Модульное и интеграционное тестирование backend}

Наиболее развито тестирование в \texttt{mobile-backend}. В репозитории присутствуют тесты для:

\begin{itemize}[leftmargin=1.2cm]
    \item регистрации пользователя;
    \item получения рыночной котировки;
    \item чтения деталей владения акцией;
    \item расчёта статистики портфеля;
    \item логики продажи акций;
    \item репозитория пользователей;
    \item HTTP-маршрутов аутентификации;
    \item HTTP-маршрутов рыночных данных.
\end{itemize}

Для \texttt{quotes-service} присутствуют тесты парсинга, store и HTTP-обработчиков. Следовательно, серверная часть уже имеет реальную основу автоматизированной проверки, а не только декларативное описание.

\subsubsection{Недостающие области}

У \texttt{mobile-api} подключены зависимости для тестирования, но собственного каталога \texttt{src/test} пока нет. У \texttt{android-app} присутствуют шаблонные зависимости для unit и instrumentation testing, однако содержательных тестов пока не реализовано.

\subsection{Предлагаемая стратегия дальнейшего развития}

\subsubsection{Модульное тестирование}

Должно покрывать:

\begin{itemize}[leftmargin=1.2cm]
    \item use case backend;
    \item парсинг и преобразование котировок;
    \item DTO и data layer Android-клиента после интеграции с API;
    \item вспомогательные функции observability и request context.
\end{itemize}

\subsubsection{Интеграционное тестирование}

Должно покрывать:

\begin{itemize}[leftmargin=1.2cm]
    \item \texttt{mobile-api} + mock upstream;
    \item \texttt{mobile-backend} + тестовая БД;
    \item \texttt{quotes-service} + тестовые Redis/ClickHouse doubles;
    \item Android data layer + mock web server.
\end{itemize}

\subsubsection{Системное тестирование}

Минимальный end-to-end сценарий должен включать:

\begin{enumerate}[leftmargin=1.2cm]
    \item регистрацию нового пользователя;
    \item вход и получение JWT;
    \item запрос актуальной котировки;
    \item выполнение операции покупки;
    \item чтение статистики портфеля;
    \item проверку появления trace в observability stack.
\end{enumerate}

\subsection{Тестовое инструментальное обеспечение}

Для системных и интеграционных сценариев полезно разработать небольшое инструментальное ПО на Kotlin, Go или Python, как это и предлагается в материалах курса~[1]. Практически это может быть:

\begin{itemize}[leftmargin=1.2cm]
    \item CLI-клиент или набор сценариев для вызова API;
    \item нагрузочный генератор пользовательских сессий;
    \item набор smoke-сценариев для проверки Docker-стенда;
    \item утилита проверки наличия trace/span после сценария сделки.
\end{itemize}

\subsection{Вывод по тестированию}

Система тестирования уже имеет работающую основу в backend- и Go-частях, но в целом ещё не завершена. Для полного соответствия архитектуре проекта необходимо выровнять тестовую зрелость \texttt{mobile-api} и Android-клиента и добавить сквозные системные сценарии.

\section*{Заключение}
\addcontentsline{toc}{section}{Заключение}

В ходе работы был заново собран комплект отчётных материалов по актуальному состоянию проекта после обновления общего репозитория и импорта Android-клиента. Подготовлено техническое задание, описана реальная архитектура системы, зафиксированы особенности реализации и предложена структура системы тестирования.

Главный результат состоит в том, что проект уже можно рассматривать как многокомпонентную мобильную платформу с выделенным API gateway, транзакционным backend, сервисом котировок, отдельным инфраструктурным контуром и встроенной observability-подсистемой. При этом также честно зафиксированы ограничения: Android-клиент пока находится на стадии интеграции с реальным backend, а тестовая зрелость разных модулей остаётся неоднородной.

Подготовленные документация, база знаний и LaTeX-отчёт создают основу не только для сдачи по курсу, но и для дальнейшего развития проекта внутри команды.

\begin{thebibliography}{99}

\bibitem{lecture}
Ключев А.\,О. Разработка мобильных приложений. Лекция 1: введение и практическое задание. Учебные материалы курса, 2026.

\bibitem{cleanarch}
Martin R.\,C. Clean Architecture: A Craftsman's Guide to Software Structure and Design. Boston: Prentice Hall, 2017.

\bibitem{ktor}
JetBrains. Ktor Documentation [Электронный ресурс]. URL: \url{https://ktor.io/docs/server-create-and-configure.html} (дата обращения: 01.04.2026).

\bibitem{android}
Android Developers. Activities and Fragments [Электронный ресурс]. URL: \url{https://developer.android.com/guide/fragments} (дата обращения: 01.04.2026).

\bibitem{exposed}
JetBrains. Exposed Documentation [Электронный ресурс]. URL: \url{https://www.jetbrains.com/help/exposed/about.html} (дата обращения: 01.04.2026).

\bibitem{otel}
OpenTelemetry. Collector and Tracing Documentation [Электронный ресурс]. URL: \url{https://opentelemetry.io/docs/collector/} (дата обращения: 01.04.2026).

\bibitem{redis}
Redis. Streams [Электронный ресурс]. URL: \url{https://redis.io/docs/latest/develop/data-types/streams/} (дата обращения: 01.04.2026).

\bibitem{clickhouse}
ClickHouse. HTTP Interface Documentation [Электронный ресурс]. URL: \url{https://clickhouse.com/docs/en/interfaces/http} (дата обращения: 01.04.2026).

\bibitem{repo-docs}
Документация проекта \texttt{stock-tracker-app}: \texttt{README.md}, \texttt{docs/system-architecture.md}, \texttt{docs/backend-implementation.md}, \texttt{docs/endpoints.md}. Репозиторий команды, состояние на 01.04.2026.

\end{thebibliography}

\end{document}
